\documentclass[a4paper, 11pt]{article}
\usepackage[utf8]{inputenc}
\usepackage[T1]{fontenc}
\usepackage[italian]{babel}
\usepackage{geometry}
\usepackage{graphicx}
\usepackage{tabularx}
\usepackage{fancyhdr}
\usepackage{hyperref}

% Margini
\geometry{top=3cm, bottom=3cm, left=2.5cm, right=2.5cm}

% Intestazione e piè di pagina
\pagestyle{fancy}
\fancyhf{}
\fancyhead[L]{\small \textit{0x06tudio} -- Verbale di Riunione}
\fancyhead[R]{\small Anno Accademico 2026/2027}
\fancyfoot[C]{\thepage}
\renewcommand{\headrulewidth}{0.4pt}

\begin{document}
\thispagestyle{plain}

\begin{minipage}{0.6\textwidth}
    \vspace{-2cm}
    {\huge \textbf{0x06tudio}} \\[0.5em]
    {\large Progetto didattico di Ingegneria del Software} \\[0.3em]
    {\normalsize Università degli Studi di Padova}
\end{minipage}
\begin{minipage}{0.4\textwidth}
    \raggedleft
    \vspace{0pt}
    \includegraphics[width=6cm]{../../media/logo.png}
    \vspace{1cm}
    \small
    \href{mailto:0x06tudio@gmail.com}{0x06tudio@gmail.com}
\end{minipage}

\vspace{0.5cm}

\begin{center}
    {\LARGE \textbf{Verbale Riunione}}
\end{center}

\vspace{0.5cm}


\begin{tabularx}{\textwidth}{l X}
    \textbf{Versione:} & 0.1.3 \\
    \textbf{Data:} & 2026/10/22 \\
    \textbf{Uso:} & Esterno \\
    \textbf{Partecipanti (Presenti):} &  Nicolò Beltrame, Pietro Brocadello, Alberto Cavalletto, Catalin Chirilov, Davide Colabove, Andrea Maggio\\
    \textbf{Destinatari:} & Vimar SpA; Proff. Vardanega, Cardin \\
\end{tabularx}

\vspace{1cm}

\section*{Registro delle modifiche}
\begin{tabularx}{\textwidth}{|c|c|X|c|}
    \hline
    \textbf{Versione} & \textbf{Data} & \textbf{Descrizione} & \textbf{Autore} \\ \hline
    0.1.3 & 2026/10/28 & Revisione e correzione documento & Pietro Brocadello \\ \hline
    0.1.2 & 2026/10/24 & Correzione hyperref al logo di gruppo & Davide Colabove \\ \hline
    0.1.1 & 2026/10/24 & Aggiunta sezione risorse e strumenti & Pietro Brocadello \\ \hline
    0.1.0 & 2026/10/23 & Creazione e stesura bozza iniziale & Andrea Maggio \\ \hline
\end{tabularx}

\newpage
\tableofcontents
\newpage


\section{Informazioni generali}
\begin{tabularx}{\textwidth}{l X}
    \textbf{Modalità / Luogo:} & Online / Microsoft Teams \\
    \textbf{Data svolgimento:} & 2026/10/22 \\
    \textbf{Ora inizio:} & 16:00 \\
    \textbf{Ora fine:} & 17:00 \\
    \textbf{Durata totale:} & 1 ora \\
    \textbf{Scopo della riunione:} & Domande riguardanti AURA (capitolato C1) \\
    \textbf{Partecipanti esterni:} & Mariano Sciacco, Domenico Bevilacqua (Vimar SpA)
\end{tabularx}

\newpage


\section{Ordine del giorno}
Durante la riunione abbiamo affrontato i seguenti punti:
\begin{enumerate}
    \item Presentazione temi emersi dal sondaggio dell'azienda e loro risposte
    \item Live Q&A con i referenti aziendali
\end{enumerate}

\newpage


\section{Resoconto della discussione}

\subsection{Architettura e linguaggi:}
\begin{itemize}
    \item Libertà sul linguaggio da usare per l'MCP, consigliato Python + FastMCP. \vspace{-0.1cm}
    \item Il client MCP da sviluppare serve principalmente per test di unità, integrazione e verifica del funzionamento del server MCP  \vspace{-0.1cm}
    \item Va sviluppata una soluzione principalmente back-end che funzioni su servizi AI richesti dal capitolato. \vspace{-0.1cm}
\end{itemize}
\subsection{Protocolli e sicurezza:}
\begin{itemize}
    \item Sarà necessario un account MyVimar con accesso all'impianto (che verrà fornito alla consegna del kit): l'MCP dovrà associare ogni utente alla sua autorizzazione OAuth2. \vspace{-0.1cm}
    \item Per bloccare richieste sbagliate si possono configurare permessi di accesso (sola lettura o lettura e scrittura) diversi a dispositivi diversi. \vspace{-0.1cm}
    \item Per lo stato dei dispositivi e le notifiche è accettabile il polling verso API. Altre soluzioni esplorabili sono trigger via webhook. \vspace{-0.1cm}
    \item L'assistente deve saper gestire le richieste: ci possono essere varianti, sinonimi e richieste ambigue. I casi basi e avanzati presentati dall'azienda vanno presi come punto di partenza noi dovremmo estenderli. 
          Le interazioni gestite che non sono elencate nel capitolato sono \textit{nice-to-have}. \vspace{-0.1cm}
\end{itemize}
\subsection{Sviluppo e test:}
\begin{itemize}
    \item Il kit di impianto si condivide via cloud: è sufficiente che un membro del gruppo lo tenga acceso e dal suo account amministratore MyVimar o anche altri account utente che vengono aggiunti all'impianto possono permettere l'accesso agli altri.\vspace{-0.1cm}
    \item Per sviluppare in parallelo si può generare un mock server in locale. \vspace{-0.1cm}
    \item Nei test automatici si può usare un mock del servizio AI per avere test ripetibili e deterministici. \vspace{-0.1cm}
\end{itemize}
\subsection{Organizzazione:}
\begin{itemize}
    \item Per il PoC è consigliato avere un MCP funzionante su pochi casi base. Per l'MVP i requisiti nice-to-have sono preferibili agli opzionali. \vspace{-0.1cm}
    \item SAL su Teams e disponibilità di approfondimenti mirati con il team previa richiesta con adeguato anticipo. \vspace{-0.1cm}
\end{itemize}
\subsection{Risorse e strumenti:}
\begin{itemize}
    \item I referenti aziendali hanno risposto che, qualora fosse necessario, previa opportuna motivazione e richiesta da parte del gruppo, forniranno le risorse menzionate nel capitolato senza problemi. \vspace{-0.1cm}
\end{itemize}
\newpage


\section{Approvazione Esterna}

Il presente verbale è stato condiviso con il referente aziendale in data 2026/10/27 e successivamente revisionato e reinviato in data 2026/10/28 ed è da ritenersi approvato e conforme a quanto discusso durante il colloquio. 

\vspace{3cm}

\end{document}